% =====================================================================
%  CHAPTER 99: 英文缩写与专有名词汇总
% =====================================================================
\chapter{英文缩写与专有名词汇总}
\chapterintro{
本附录收录《营养与食品卫生学》课程中出现的全部英文缩写及专有名词，按类别分为九大类，以表格形式呈现，方便对照复习和查阅。
}

% ===== 一、国际组织与机构 =====
\section{一、国际组织与机构}

\begin{table}[H]
\centering
\begin{tabularx}{\textwidth}{lX}
\toprule
\textbf{缩写} & \textbf{全称与中文释义} \\
\midrule
WHO & World Health Organization \hfill 世界卫生组织 \\
FAO & Food and Agriculture Organization \hfill 联合国粮食及农业组织 \\
CAC & Codex Alimentarius Commission \hfill 国际食品法典委员会 \\
JECFA & Joint FAO/WHO Expert Committee on Food Additives \hfill FAO/WHO食品添加剂联合专家委员会 \\
IARC & International Agency for Research on Cancer \hfill 国际癌症研究中心 \\
EFSA & European Food Safety Authority \hfill 欧洲食品安全局 \\
WCRF & World Cancer Research Fund \hfill 世界癌症研究基金会 \\
AICR & American Institute for Cancer Research \hfill 美国癌症研究所 \\
WTO & World Trade Organization \hfill 世界贸易组织 \\
ISO & International Organization for Standardization \hfill 国际标准化组织 \\
\bottomrule
\end{tabularx}
\end{table}

% ===== 二、营养学基础 =====
\section{二、营养学基础}

\subsection*{（一）氨基酸与蛋白质}

\begin{table}[H]
\centering
\begin{tabularx}{\textwidth}{lX}
\toprule
\textbf{缩写} & \textbf{全称与中文释义} \\
\midrule
EAA & Essential Amino Acid \hfill 必需氨基酸 \\
Ile & Isoleucine \hfill 异亮氨酸 \\
Leu & Leucine \hfill 亮氨酸 \\
Lys & Lysine \hfill 赖氨酸 \\
Met & Methionine \hfill 蛋氨酸（甲硫氨酸） \\
Phe & Phenylalanine \hfill 苯丙氨酸 \\
Thr & Threonine \hfill 苏氨酸 \\
Trp & Tryptophan \hfill 色氨酸 \\
Val & Valine \hfill 缬氨酸 \\
His & Histidine \hfill 组氨酸（婴儿必需氨基酸） \\
Cys & Cysteine \hfill 半胱氨酸（条件必需氨基酸） \\
Tyr & Tyrosine \hfill 酪氨酸（条件必需氨基酸） \\
ONL & Obligatory Nitrogen Loss \hfill 必要氮损失 \\
PEM & Protein-Energy Malnutrition \hfill 蛋白质-能量营养不良 \\
AD & Apparent Digestibility \hfill 表观消化率 \\
TD & True Digestibility \hfill 真消化率 \\
BV & Biological Value \hfill 生物价 \\
NPU & Net Protein Utilization \hfill 蛋白质净利用率 \\
PER & Protein Efficiency Ratio \hfill 蛋白质功效比值 \\
AAS & Amino Acid Score \hfill 氨基酸评分 \\
PDCAAS & Protein Digestibility Corrected Amino Acid Score \hfill 经消化率修正的氨基酸评分 \\
DIAAS & Digestible Indispensable Amino Acid Score \hfill 可消化必需氨基酸评分 \\
LAA & Limiting Amino Acid \hfill 限制氨基酸 \\
\bottomrule
\end{tabularx}
\end{table}

\subsection*{（二）脂类}

\begin{table}[H]
\centering
\begin{tabularx}{\textwidth}{lX}
\toprule
\textbf{缩写} & \textbf{全称与中文释义} \\
\midrule
SFA & Saturated Fatty Acid \hfill 饱和脂肪酸 \\
MUFA & Monounsaturated Fatty Acid \hfill 单不饱和脂肪酸 \\
PUFA & Polyunsaturated Fatty Acid \hfill 多不饱和脂肪酸 \\
EFA & Essential Fatty Acid \hfill 必需脂肪酸 \\
DHA & Docosahexaenoic Acid \hfill 二十二碳六烯酸 \\
EPA & Eicosapentaenoic Acid \hfill 二十碳五烯酸 \\
AA & Arachidonic Acid \hfill 花生四烯酸 \\
ARA & Arachidonic Acid \hfill 花生四烯酸（同AA） \\
LA & Linoleic Acid \hfill 亚油酸 \\
ALA & $\alpha$-Linolenic Acid \hfill $\alpha$-亚麻酸 \\
MCFA & Medium Chain Fatty Acid \hfill 中链脂肪酸 \\
SCFA & Short Chain Fatty Acid \hfill 短链脂肪酸 \\
TFA & Trans Fatty Acid \hfill 反式脂肪酸 \\
LCT & Long Chain Triglyceride \hfill 长链甘油三酯 \\
MCT & Medium Chain Triglyceride \hfill 中链甘油三酯 \\
\bottomrule
\end{tabularx}
\end{table}

\subsection*{（三）碳水合物与能量}

\begin{table}[H]
\centering
\begin{tabularx}{\textwidth}{lX}
\toprule
\textbf{缩写} & \textbf{全称与中文释义} \\
\midrule
GI & Glycemic Index \hfill 血糖生成指数 \\
GL & Glycemic Load \hfill 血糖负荷 \\
BEE & Basal Energy Expenditure \hfill 基础代谢能量消耗 \\
BMR & Basal Metabolic Rate \hfill 基础代谢率 \\
PAL & Physical Activity Level \hfill 体力活动水平 \\
TEF & Thermic Effect of Food \hfill 食物热效应 \\
SDA & Specific Dynamic Action \hfill 食物特殊动力作用（即TEF） \\
RQ & Respiratory Quotient \hfill 呼吸商 \\
REE & Resting Energy Expenditure \hfill 静息能量消耗 \\
\bottomrule
\end{tabularx}
\end{table}

\subsection*{（四）维生素与辅酶}

\begin{table}[H]
\centering
\begin{tabularx}{\textwidth}{lX}
\toprule
\textbf{缩写} & \textbf{全称与中文释义} \\
\midrule
RE & Retinol Equivalent \hfill 视黄醇当量 \\
RAE & Retinol Activity Equivalent \hfill 视黄醇活性当量 \\
TPP & Thiamine Pyrophosphate \hfill 硫胺素焦磷酸酯（维生素B1活性形式） \\
TDP & Thiamine Diphosphate \hfill 硫胺素二磷酸（同TPP） \\
FMN & Flavin Mononucleotide \hfill 黄素单核苷酸 \\
FAD & Flavin Adenine Dinucleotide \hfill 黄素腺嘌呤二核苷酸 \\
NAD$^+$ & Nicotinamide Adenine Dinucleotide \hfill 烟酰胺腺嘌呤二核苷酸（辅酶I） \\
NADP$^+$ & Nicotinamide Adenine Dinucleotide Phosphate \hfill 烟酰胺腺嘌呤二核苷酸磷酸（辅酶II） \\
THF & Tetrahydrofolate \hfill 四氢叶酸 \\
CoA & Coenzyme A \hfill 辅酶A \\
ACP & Acyl Carrier Protein \hfill 酰基载体蛋白 \\
NE & Niacin Equivalent \hfill 烟酸当量 \\
DFE & Dietary Folate Equivalent \hfill 膳食叶酸当量 \\
$\alpha$-TE & $\alpha$-Tocopherol Equivalent \hfill $\alpha$-生育酚当量 \\
PLP & Pyridoxal Phosphate \hfill 磷酸吡哆醛（维生素B6活性形式） \\
MTHF & 5-Methyltetrahydrofolate \hfill 5-甲基四氢叶酸 \\
\bottomrule
\end{tabularx}
\end{table}

\subsection*{（五）矿物质}

\begin{table}[H]
\centering
\begin{tabularx}{\textwidth}{lX}
\toprule
\textbf{缩写} & \textbf{全称与中文释义} \\
\midrule
Ca & Calcium \hfill 钙 \\
Fe & Iron \hfill 铁 \\
Zn & Zinc \hfill 锌 \\
Se & Selenium \hfill 硒 \\
I & Iodine \hfill 碘 \\
Mg & Magnesium \hfill 镁 \\
Cu & Copper \hfill 铜 \\
Cr & Chromium \hfill 铬 \\
Mn & Manganese \hfill 锰 \\
F & Fluorine \hfill 氟 \\
P & Phosphorus \hfill 磷 \\
K & Potassium \hfill 钾 \\
Na & Sodium \hfill 钠 \\
Cl & Chlorine \hfill 氯 \\
\bottomrule
\end{tabularx}
\end{table}

% ===== 三、公共营养 =====
\section{三、公共营养}

\begin{table}[H]
\centering
\begin{tabularx}{\textwidth}{lX}
\toprule
\textbf{缩写} & \textbf{全称与中文释义} \\
\midrule
DRIs & Dietary Reference Intakes \hfill 膳食营养素参考摄入量 \\
EAR & Estimated Average Requirement \hfill 平均需要量 \\
RNI & Recommended Nutrient Intake \hfill 推荐摄入量 \\
AI & Adequate Intake \hfill 适宜摄入量 \\
UL & Tolerable Upper Intake Level \hfill 可耐受最高摄入量 \\
AMDR & Acceptable Macronutrient Distribution Range \hfill 宏量营养素可接受范围 \\
PI-NCD & Proposed Intake for Non-communicable Chronic Disease \hfill 预防非传染病慢性病的建议摄入量 \\
SPL & Specific Proposed Level \hfill 特定建议值 \\
BMI & Body Mass Index \hfill 体质指数 \\
NRV & Nutrient Reference Value \hfill 营养素参考值 \\
FFQ & Food Frequency Questionnaire \hfill 食物频率法 \\
NRS-2002 & Nutrition Risk Screening 2002 \hfill 营养风险筛查2002 \\
SGA & Subjective Global Assessment \hfill 主观全面评定 \\
MNA & Mini Nutritional Assessment \hfill 微型营养评定 \\
MUST & Malnutrition Universal Screening Tool \hfill 营养不良通用筛查工具 \\
DDP & Desirable Dietary Pattern \hfill 理想膳食模式 \\
DGE & German Nutrition Society \hfill 德国营养学会（DDP模型来源） \\
\bottomrule
\end{tabularx}
\end{table}

% ===== 四、临床营养 =====
\section{四、临床营养}

\begin{table}[H]
\centering
\begin{tabularx}{\textwidth}{lX}
\toprule
\textbf{缩写} & \textbf{全称与中文释义} \\
\midrule
EN & Enteral Nutrition \hfill 肠内营养 \\
PN & Parenteral Nutrition \hfill 肠外营养 \\
TPN & Total Parenteral Nutrition \hfill 全肠外营养 \\
PPN & Peripheral Parenteral Nutrition \hfill 周围静脉营养 \\
FSMP & Food for Special Medical Purposes \hfill 特殊医学用途配方食品 \\
DASH & Dietary Approaches to Stop Hypertension \hfill 终止高血压膳食疗法 \\
GDM & Gestational Diabetes Mellitus \hfill 妊娠期糖尿病 \\
DOHaD & Developmental Origins of Health and Disease \hfill 健康和疾病的发育起源（都哈理论） \\
T2DM & Type 2 Diabetes Mellitus \hfill 2型糖尿病 \\
CVD & Cardiovascular Disease \hfill 心血管疾病 \\
CHD & Coronary Heart Disease \hfill 冠心病 \\
NCD & Non-communicable Chronic Disease \hfill 非传染性慢性病 \\
\bottomrule
\end{tabularx}
\end{table}

% ===== 五、食品污染 =====
\section{五、食品污染}

\subsection*{（一）微生物与腐败}

\begin{table}[H]
\centering
\begin{tabularx}{\textwidth}{lX}
\toprule
\textbf{缩写} & \textbf{全称与中文释义} \\
\midrule
CFU & Colony Forming Unit \hfill 菌落形成单位 \\
MPN & Most Probable Number \hfill 大肠菌群最大可能数 \\
Aw & Water Activity \hfill 水分活性 \\
TVBN & Total Volatile Basic Nitrogen \hfill 挥发性盐基总氮 \\
TMA & Trimethylamine \hfill 三甲胺 \\
D值 & Decimal Reduction Time \hfill 十进制减菌时间（90\%杀灭时间） \\
F值 & Sterilization Time \hfill 杀菌值 \\
Z值 & Temperature Change for 10-fold D-value Change \hfill Z值 \\
UHT & Ultra High Temperature \hfill 超高温消毒法 \\
LTLT & Low Temperature Long Time \hfill 低温长时间巴氏消毒 \\
HTST & High Temperature Short Time \hfill 高温短时间巴氏消毒 \\
TTT & Time-Temperature-Tolerance \hfill 时间-温度-耐受原则 \\
MAP & Modified Atmosphere Packaging \hfill 气调包装 \\
\bottomrule
\end{tabularx}
\end{table}

\subsection*{（二）化学污染物与农药}

\begin{table}[H]
\centering
\begin{tabularx}{\textwidth}{lX}
\toprule
\textbf{缩写} & \textbf{全称与中文释义} \\
\midrule
MRLs & Maximum Residue Limits \hfill 最大残留限量 \\
ADI & Acceptable Daily Intake \hfill 每日允许摄入量 \\
BHA & Butylated Hydroxyanisole \hfill 丁基羟基茴香醚（抗氧化剂） \\
BHT & Butylated Hydroxytoluene \hfill 二丁基羟基甲苯（抗氧化剂） \\
PG & Propyl Gallate \hfill 没食子酸丙酯（抗氧化剂） \\
TBHQ & Tertiary Butylhydroquinone \hfill 特丁基对苯二酚（抗氧化剂） \\
DDT & Dichlorodiphenyltrichloroethane \hfill 二氯二苯三氯乙烷（有机氯农药） \\
BHC/666 & Benzene Hexachloride \hfill 六六六（有机氯农药） \\
PAH & Polycyclic Aromatic Hydrocarbon \hfill 多环芳烃 \\
B(a)P & Benzo(a)pyrene \hfill 苯并(a)芘 \\
NDMA & N-Nitrosodimethylamine \hfill 二甲基亚硝胺 \\
NDEA & N-Nitrosodiethylamine \hfill 二乙基亚硝胺 \\
3-MCPD & 3-Monochloropropane-1,2-diol \hfill 3-氯-1,2-丙二醇 \\
PCB & Polychlorinated Biphenyl \hfill 多氯联苯 \\
POPs & Persistent Organic Pollutants \hfill 持久性有机污染物 \\
\bottomrule
\end{tabularx}
\end{table}

\subsection*{（三）真菌毒素}

\begin{table}[H]
\centering
\begin{tabularx}{\textwidth}{lX}
\toprule
\textbf{缩写} & \textbf{全称与中文释义} \\
\midrule
AF & Aflatoxin \hfill 黄曲霉毒素 \\
AFB1 & Aflatoxin B1 \hfill 黄曲霉毒素B1（毒性和致癌性最强） \\
AFM1 & Aflatoxin M1 \hfill 黄曲霉毒素M1（乳中代谢产物） \\
OTA & Ochratoxin A \hfill 赭曲霉毒素A \\
DON & Deoxynivalenol \hfill 脱氧雪腐镰刀菌烯醇（呕吐毒素） \\
3-NPA & 3-Nitropropionic Acid \hfill 3-硝基丙酸（霉变甘蔗毒素） \\
ZEN/ZEA & Zearalenone \hfill 玉米赤霉烯酮 \\
T-2 & T-2 Toxin \hfill T-2毒素 \\
ATA & Alimentary Toxic Aleukia \hfill 食物中毒性白细胞缺乏症（T-2引起） \\
NIV & Nivalenol \hfill 雪腐镰刀菌烯醇 \\
\bottomrule
\end{tabularx}
\end{table}

\subsection*{（四）ATP分解产物}

\begin{table}[H]
\centering
\begin{tabularx}{\textwidth}{lX}
\toprule
\textbf{缩写} & \textbf{全称与中文释义} \\
\midrule
ATP & Adenosine Triphosphate \hfill 腺苷三磷酸 \\
ADP & Adenosine Diphosphate \hfill 腺苷二磷酸 \\
AMP & Adenosine Monophosphate \hfill 腺苷一磷酸 \\
IMP & Inosine Monophosphate \hfill 肌苷一磷酸（鲜味物质） \\
HxR & Inosine \hfill 肌苷/次黄嘌呤核苷 \\
Hx & Hypoxanthine \hfill 次黄嘌呤 \\
K值 & Freshness Index K value \hfill 鱼肉鲜度指标K值 \\
\bottomrule
\end{tabularx}
\end{table}

% ===== 六、食品安全风险分析与监管 =====
\section{六、食品安全风险分析与监管}

\begin{table}[H]
\centering
\begin{tabularx}{\textwidth}{lX}
\toprule
\textbf{缩写} & \textbf{全称与中文释义} \\
\midrule
HBGV & Health-Based Guidance Value \hfill 健康指导值 \\
POD & Point of Departure \hfill 出发点 \\
NOAEL & No Observed Adverse Effect Level \hfill 未观察到有害作用水平 \\
LOAEL & Lowest Observed Adverse Effect Level \hfill 观察到有害作用最低水平 \\
BMDL & Benchmark Dose Lower Confidence Limit \hfill 基准剂量下限 \\
UFs & Uncertainty Factors \hfill 不确定系数 \\
ARfD & Acute Reference Dose \hfill 急性参考剂量 \\
TDI & Tolerable Daily Intake \hfill 日耐受摄入量 \\
PTWI & Provisional Tolerable Weekly Intake \hfill 暂定每周耐受摄入量 \\
PTMI & Provisional Tolerable Monthly Intake \hfill 暂定每月耐受摄入量 \\
MOE & Margin of Exposure \hfill 暴露边界 \\
ALARA & As Low As Reasonably Achievable \hfill 合理可行的最低水平 \\
GMP & Good Manufacturing Practice \hfill 食品良好生产规范 \\
SSOP & Sanitation Standard Operating Procedure \hfill 卫生标准操作程序 \\
HACCP & Hazard Analysis and Critical Control Point \hfill 危害分析及关键控制点 \\
CCP & Critical Control Point \hfill 关键控制点 \\
CL & Critical Limit \hfill 关键限值 \\
OL & Operating Limit \hfill 操作限值 \\
HA & Hazard Analysis \hfill 危害分析 \\
CA & Corrective Action \hfill 纠偏程序 \\
V & Verification \hfill 验证程序 \\
R & Record-keeping \hfill 记录保存 \\
GRAS & Generally Recognized As Safe \hfill 公认安全 \\
SAMR & State Administration for Market Regulation \hfill 国家市场监督管理总局 \\
CFDA & China Food and Drug Administration \hfill 国家食品药品监督管理总局（前身） \\
GHP & Good Hygiene Practice \hfill 良好卫生规范 \\
TDS & Total Diet Study \hfill 总膳食研究 \\
\bottomrule
\end{tabularx}
\end{table}

% ===== 七、食源性疾病 =====
\section{七、食源性疾病}

\begin{table}[H]
\centering
\begin{tabularx}{\textwidth}{lX}
\toprule
\textbf{缩写} & \textbf{全称与中文释义} \\
\midrule
BSE & Bovine Spongiform Encephalopathy \hfill 疯牛病 \\
nvCJD & new variant Creutzfeldt-Jakob Disease \hfill 新型变异型克雅氏病 \\
PrP$^\text{C}$ & Cellular Prion Protein \hfill 细胞型朊蛋白（正常型） \\
PrP$^\text{Sc}$ & Scrapie Prion Protein \hfill 致病型朊病毒蛋白 \\
TTX & Tetrodotoxin \hfill 河豚毒素 \\
HCN & Hydrogen Cyanide \hfill 氢氰酸 \\
MNL & Maximum No-effect Level \hfill 最大无作用剂量 \\
3-NPA & 3-Nitropropionic Acid \hfill 3-硝基丙酸（霉变甘蔗） \\
FMD & Foot and Mouth Disease \hfill 口蹄疫 \\
HUS & Hemolytic Uremic Syndrome \hfill 溶血性尿毒综合征（E. coli O157引起） \\
ETEC & Enterotoxigenic E. coli \hfill 肠产毒性大肠杆菌 \\
EHEC & Enterohemorrhagic E. coli \hfill 肠出血性大肠杆菌 \\
VTEC & Vero-toxin producing E. coli \hfill 产Vero毒素大肠杆菌 \\
STEC & Shiga-toxin producing E. coli \hfill 产志贺毒素大肠杆菌 \\
\bottomrule
\end{tabularx}
\end{table}

% ===== 八、其他食品营养相关概念 =====
\section{八、食品营养相关其他概念}

\begin{table}[H]
\centering
\begin{tabularx}{\textwidth}{lX}
\toprule
\textbf{缩写} & \textbf{全称与中文释义} \\
\midrule
GMO & Genetically Modified Organism \hfill 转基因生物/转基因食品 \\
GSH-Px & Glutathione Peroxidase \hfill 谷胱甘肽过氧化物酶（含硒酶） \\
GTF & Glucose Tolerance Factor \hfill 葡萄糖耐量因子（铬的活性形式） \\
SOD & Superoxide Dismutase \hfill 超氧化物歧化酶 \\
IGF-1 & Insulin-like Growth Factor-1 \hfill 胰岛素样生长因子-1 \\
TNF-$\alpha$ & Tumor Necrosis Factor-$\alpha$ \hfill 肿瘤坏死因子-$\alpha$ \\
IL-6 & Interleukin-6 \hfill 白细胞介素-6 \\
CRP & C-Reactive Protein \hfill C反应蛋白 \\
HCG & Human Chorionic Gonadotropin \hfill 人绒毛膜促性腺激素 \\
HCS & Human Chorionic Somatomammotropin \hfill 人绒毛膜生长催乳素 \\
TSH & Thyroid Stimulating Hormone \hfill 促甲状腺激素 \\
T3 & Triiodothyronine \hfill 三碘甲状腺原氨酸 \\
T4 & Thyroxine \hfill 甲状腺素 \\
Hb & Hemoglobin \hfill 血红蛋白 \\
SF & Serum Ferritin \hfill 血清铁蛋白 \\
sTfR & Soluble Transferrin Receptor \hfill 血清运铁蛋白受体 \\
FEP & Free Erythrocyte Protoporphyrin \hfill 红细胞游离原卟啉 \\
TIBC & Total Iron Binding Capacity \hfill 总铁结合力 \\
LDL & Low Density Lipoprotein \hfill 低密度脂蛋白 \\
HDL & High Density Lipoprotein \hfill 高密度脂蛋白 \\
VLDL & Very Low Density Lipoprotein \hfill 极低密度脂蛋白 \\
IDL & Intermediate Density Lipoprotein \hfill 中间密度脂蛋白 \\
TC & Total Cholesterol \hfill 总胆固醇 \\
TG & Triglycerides \hfill 甘油三酯 \\
LDL-C & LDL Cholesterol \hfill 低密度脂蛋白胆固醇 \\
HDL-C & HDL Cholesterol \hfill 高密度脂蛋白胆固醇 \\
FFA & Free Fatty Acid \hfill 游离脂肪酸 \\
BP & Blood Pressure \hfill 血压 \\
SBP & Systolic Blood Pressure \hfill 收缩压 \\
DBP & Diastolic Blood Pressure \hfill 舒张压 \\
OGTT & Oral Glucose Tolerance Test \hfill 口服葡萄糖耐量试验 \\
HbA1c & Glycated Hemoglobin \hfill 糖化血红蛋白 \\
RDA & Recommended Dietary Allowance \hfill 推荐膳食供给量（旧DRIs体系） \\
RNI & Recommended Nutrient Intake \hfill 推荐摄入量（新DRIs体系） \\
\bottomrule
\end{tabularx}
\end{table}

% ===== 九、氨基酸三字母/单字母对照 =====
\section{九、常见氨基酸三字母与单字母缩写对照}

\begin{table}[H]
\centering
\begin{tabularx}{\textwidth}{llll}
\toprule
\textbf{中文名} & \textbf{英文名} & \textbf{三字母} & \textbf{单字母} \\
\midrule
甘氨酸 & Glycine & Gly & G \\
丙氨酸 & Alanine & Ala & A \\
缬氨酸* & Valine & Val & V \\
亮氨酸* & Leucine & Leu & L \\
异亮氨酸* & Isoleucine & Ile & I \\
苯丙氨酸* & Phenylalanine & Phe & F \\
色氨酸* & Tryptophan & Trp & W \\
蛋氨酸* & Methionine & Met & M \\
脯氨酸 & Proline & Pro & P \\
丝氨酸 & Serine & Ser & S \\
苏氨酸* & Threonine & Thr & T \\
半胱氨酸 & Cysteine & Cys & C \\
酪氨酸 & Tyrosine & Tyr & Y \\
天冬酰胺 & Asparagine & Asn & N \\
谷氨酰胺 & Glutamine & Gln & Q \\
天冬氨酸 & Aspartic acid & Asp & D \\
谷氨酸 & Glutamic acid & Glu & E \\
赖氨酸* & Lysine & Lys & K \\
精氨酸 & Arginine & Arg & R \\
组氨酸（婴儿*） & Histidine & His & H \\
\bottomrule
\multicolumn{4}{l}{\small * 标注为必需氨基酸（成人8种+婴儿组氨酸）} \\
\end{tabularx}
\end{table}

% ===== 十、常用希腊字母与符号 =====
\section{十、常用希腊字母与符号}

\begin{table}[H]
\centering
\begin{tabularx}{\textwidth}{llX}
\toprule
\textbf{符号} & \textbf{名称} & \textbf{常见用途} \\
\midrule
$\alpha$ & Alpha & $\alpha$-亚麻酸、$\alpha$-生育酚、$\alpha$-螺旋 \\
$\beta$ & Beta & $\beta$-胡萝卜素、$\beta$-折叠、$\beta$-氧化 \\
$\gamma$ & Gamma & $\gamma$-生育酚、$\gamma$-亚麻酸 \\
$\delta$ & Delta & $\delta$-生育酚 \\
$\Delta$ & Delta（大写） & 表示双键位置（如$\Delta$9-desaturase） \\
$\omega$ & Omega & $\omega$-3 / $\omega$-6脂肪酸 \\
$\mu$ & Mu & 微克（$\mu$g） \\
\bottomrule
\end{tabularx}
\end{table}

\newpage
